% META: {"id":"1SPE-GEOREP-CO-013","chapitre":"1SPE-GEOMETRIE-REPEREE","type_objet":"corrige","exercice_ref":"1SPE-GEOREP-EX-013","capacites_codes":["C2"],"sources_inspiration":[],"status":"approved","origin":"generated","status_history":["generated","audited","accepted","approved"],"release_acceptance":"RELEASE_OWNER_BATCH_ACCEPTANCE","acceptance_closure_digest":"sha256:9b3ccf9a81c5520fbb7e03b7b2d3e7bf2057a02834b6d908bf3be5f49f3b553f","acceptance_receipt_digest":"sha256:4fad86f0282e0fa4e0419cca1ad6379ae7af5a280c04a4a90880f90a1c044242"}
% BEGIN-VERIFY
% from sympy import *
% # D1: 2x + 3y - 1 = 0, D2: 6x + 9y + 4 = 0
% # Vecteurs normaux: n1(2,3), n2(6,9) = 3*n1 => colineaires => paralleles
% assert 2*9 - 3*6 == 0
% END-VERIFY
\begin{corrige}{1SPE-GEOREP-EX-013}

\commentaireMarge{On compare les vecteurs normaux.}

$\vec{n}_1(2 ; 3)$ et $\vec{n}_2(6 ; 9) = 3\vec{n}_1$.

Les vecteurs normaux sont colinéaires, donc les droites sont \textbf{parallèles}.

Elles ne sont pas confondues car $\dfrac{c_2}{c_1} = \dfrac{4}{-1} = -4
\neq 3$.

\end{corrige}
